\section{Profile 23 construction}
\label{sec:construction}

Profile~23 is a transparent, WHIR-style polynomial argument over a circle
evaluation domain.  It uses the circle-code machinery developed for Circle
STARKs~\cite{HaboeckLevitPapini2024} and an arity-four folding protocol in the
lineage of FRI and WHIR~\cite{BenSassonBentovHoreshRiabzev2018FRI,ArnonChiesaFenziYogev2025WHIR}.
The transcript, oracle layout, masking system, and relation checks described
below are specific to Profile~23; the construction is not an invocation of an
unmodified WHIR parameter set.

\subsection{Transparent parameters}

\begin{definition}[Transparent parameters]
\label{def:profile23-transparent-parameters}
The public parameters \(\mathsf{pp}_{23}\) consist of the field tower,
Poseidon2 constants, circle-domain generators, fixed radix-four fold tables,
the atomic-v3 constraint and copy registries, the 10-bit lookup table, Merkle
domain tags, transcript labels, and the following shape.
\end{definition}

\begin{table}[t]
\centering
\caption{Frozen Profile~23 polynomial and transcript parameters.}
\label{tab:profile23-parameters}
\begin{tabular}{@{}ll@{}}
\toprule
quantity & value \\
\midrule
trace/message rows & \(2^{10}\) \\
circle evaluation symbols & \(2^{19}\) \\
inverse rate & \(512\) \\
layer-zero query fibers & \(N=2^{17}\) \\
queries per branch & \(q=18\), without replacement \\
generator width & \(29\) \\
generator order & semantic \(0\ldots15\), mask \(16\ldots25\), \(H_{26},G_{27},D_{28}\) \\
query branches per attempt & \(3\) \\
attempt cap & \(17\) \\
batch work & \(37\) bits \\
fold work & \(34,33,30,25\) bits \\
final work & \(32\) bits \\
folds & four arity-four folds \\
final polynomial & four coefficients in \(\mathbb K\) \\
\bottomrule
\end{tabular}
\end{table}

These parameters require no structured reference string, trapdoor, or secret
setup.  They do require the fixed hash and random-oracle assumptions made
explicit in Sections~\ref{sec:soundness} and~\ref{sec:hiding}.

\subsection{Layer-zero oracles and batching}

The layer-zero message contains 26 \(\mathbb F\)-valued codewords and three
\(\mathbb K\)-valued codewords.  The first sixteen codewords contain the
randomized semantic trace.  Ten additional full-domain codewords carry
mask-only directions.  The three extension-field lanes are denoted
\(H,G,D\).  The \(D\) lane is a full \(2^{10}\)-coefficient codeword with
zero masking factor.  It is nevertheless committed and opened as an ordinary
lane.

The first commitment, \(C_1\), stores the 26 subfield codewords.  Four circle
symbols form one 416-byte leaf.  The second commitment, \(C_2\), stores
\(H,G,D\); four extension-field symbols form one 192-byte leaf.  After the
nonzero batching challenge \(\gamma\in\mathbb K^*\), the verifier combines
the 29 lanes as
\begin{equation}
 R_\gamma=\sum_{j=0}^{25}\gamma^j C_j
              +\gamma^{26}H+\gamma^{27}G+\gamma^{28}D.
 \label{eq:profile23-batch}
\end{equation}
The zero-factor lane is designed to add hiding directions without moving the
committed batched word.  For example,
\(
  \Delta G=-\gamma\Delta D
\)
implies
\(
 \gamma^{27}\Delta G+\gamma^{28}\Delta D=0
\).
More generally, a source change \(U\) in lane \(j<27\) is paired with
\(\Delta G=-\gamma^{j-27}U\).  The nonzero sampling of \(\gamma\) makes these
expressions well defined.  The complete source family and the legal
root-neutral directions are checked by \(\mathsf{Good}_{23}\), defined below.

The circle codeword has \(2^{19}\) symbols, grouped into \(2^{17}\)
arity-four fibers.  An arity-four fold produces the first line codeword; three
more folds reduce through authenticated line layers of depths 15, 13, and 11
to a four-coefficient terminal polynomial.  Each round contains two
out-of-domain evaluations and a relation sumcheck.  At a queried layer-zero
fiber, the verifier decodes the \(C_1\) and \(C_2\) leaves, evaluates
\eqref{eq:profile23-batch}, and checks the first fold.  It then checks the
three authenticated line transitions and the terminal polynomial.

\subsection{Fiat--Shamir schedule}

The transcript uses SHA-256 with rigid, typed domain labels.  Its logical
order is as follows.
\begin{enumerate}
  \item Absorb the profile header, field-basis identifier, canonical statement
        digest, masking context, and nonzero public mask nonce.
  \item Absorb the \(C_1\) root; sample \(\lambda\) and \(\chi\).  Absorb the
        \(C_2\) root and sample the zerocheck batching challenges.
  \item Absorb the initial mask claim, sample nonzero \(\eta\), verify and
        absorb the ten-round masked sumcheck, and derive its terminal point
        \(z\in\mathbb K^{10}\).
  \item Absorb the three statement points \(z\), \(\operatorname{succ}(z)\),
        and \(\operatorname{xor}_{12}(z)\), their 84 extension-field
        evaluations, and the three \(D\)-lane evaluations.
  \item Check and absorb the 37-bit batch-work nonce; sample
        \(\gamma\in\mathbb K^*\) and the point scale.
  \item For fold round \(r=0,\ldots,3\), derive the two OOD points, absorb
        the two OOD values and relation sumcheck (and, for \(r>0\), the new
        layer root), then check and absorb the round's work nonce before
        sampling its fold challenge \(\alpha_r\).
  \item Absorb the four-coefficient final polynomial; check and absorb the
        32-bit final-work nonce.
  \item Fork the resulting state into three branches.  Branch \(b\) absorbs
        the single byte \(b\in\{0,1,2\}\) under transcript label 44 and then
        samples 18 distinct query fibers.
\end{enumerate}
The first seven steps are common to all three branches.  In particular, the
committed word, OOD checks, fold challenges, final polynomial, and all work
nonces are fixed before the branch byte is absorbed.  The three replays start
from the same pre-branch state; they are not chained.

\subsection{The complete \texorpdfstring{\(\mathsf{Good}_{23}\)}{Good23} predicate}

\begin{definition}[Public schedule]
\label{def:profile23-schedule}
A \Profile{} schedule \(s_{\rm fs}\) comprises the fixed layout identifiers,
all public transcript challenges and OOD points, fold data, and one ordered
18-tuple of distinct query fibers.  It excludes the witness, realized masks,
private salts, committed roots as opaque inputs, opened values, and retry
history.
\end{definition}

The deterministic predicate \(\mathsf{Good}_{23}(s_{\rm fs})\) reconstructs
the public linear maps induced by this schedule and accepts precisely when
all static pins, source guards, and rank conditions in
Table~\ref{tab:good23-ranks} hold.  Its source guards additionally check that
each selected root-neutral source has zero compressed terminal contribution
and zero pointwise \(\gamma\)-batched root difference.

\begin{table}[t]
\centering
\caption{Rank conditions checked by \(\mathsf{Good}_{23}\), over
\(\mathbb F\).  ``Schur'' means rank after eliminating all query coordinates.}
\label{tab:good23-ranks}
\begin{tabular}{@{}lrr@{}}
\toprule
map & required rank & target dimension \\
\midrule
root-neutral joint map & 1404 & 1404 \\
serialized terminal projection & 324 & 324 \\
terminal plus initial projection & 328 & 328 \\
\(\ker(\mathrm{initial},T_z)\) coverage & 1076 & 1076 \\
remaining \(G/D\) query map & 288 & 288 \\
remaining \(G/D\) terminal Schur map & 12 & 12 \\
inactive-balanced \(H_1\) query map & 288 & 288 \\
inactive-balanced \(H_1\) terminal Schur map & 12 & 12 \\
\bottomrule
\end{tabular}
\end{table}

The root-neutral basis is selected in increasing query-degree order:
semantic and \(D\) sources precede mask-only sources.  It contains 1068
degree-one and 336 degree-two query columns.  Its total query degree is at
most 31320.  The root-neutral \(z\)-degree is at most 41040; the two terminal
Schur blocks add 120 each, giving 41280.  The \(\gamma\)-coordinate degree is
at most 80688, and the complete continuous degree is therefore 121968.

\begin{lemma}[Complete-Good algebraic liveness]
\label{lem:good23-product}
For a uniformly sampled branch schedule, let \(B\) be the event that
\(\mathsf{Good}_{23}\) rejects.  For three independent post-final branches
from one common prefix,
\[
 \Pr[B_0\cap B_1\cap B_2]\le
 \beta:=\frac{121968}{p}+
 \left(\frac{31320}{131072-17}\right)^3
 =0.013705910239932412.
\]
With at most 17 independent, completed attempts, the probability of an
internal algebraic Abort caused by sampler/build failure or cap exhaustion is
at most
\[
  17\frac{128}{p^6}+\beta^{17}
  <2^{-105.21398677941983}.
\]
\end{lemma}

\begin{proof}
Each nonzero determinant defining the complete product has continuous degree
at most 121968, so Schwartz--Zippel bounds failure over its continuous
challenges by \(121968/p\).  The exhaustive query-domain map has
\(N=131072\) distinct roots and none equal to one.  Sampling 18 roots without
replacement therefore bounds failure of one query minor by
\(31320/(N-17)\).  The three label-44 branches are independent in the
random-oracle experiment, producing the cube.  Independent fresh attempts
raise this bound to the seventeenth power; the first term accounts for the
bounded six-coordinate sampler/build abort in each attempt.  This is an
internal algebraic availability bound, not the total probability of public
\(\mathsf{Abort}\) under a wall-clock controller.  It does not include failure
to complete an attempt before a deployment deadline.  Soundness does not
condition on the event in this lemma.
\end{proof}

\subsection{Prover and verifier}

One honest attempt proceeds in two phases.  It first builds the common masked
trace, the two layer-zero commitments, all OOD and fold data, the terminal
polynomial, and the six canonical work records.  It then evaluates
\(\mathsf{Good}_{23}\) on all three post-final schedules and selects the least
good selector.  Only the selected branch's five salted minimal-subtree
opening sections are serialized.  If all three branches are bad, the attempt
is erased and repeated with fresh entropy and a new durably burned nonce.  A
non-rank construction or schema error is fatal but has the same public abort
encoding.  After 17 unsuccessful completed attempts the private builder
records cap exhaustion.

The production executable accepts a caller-selected wall-clock boundary and
uses 3,600 seconds when none is supplied.  At the selected boundary its
controller emits a payload-free \(\mathsf{Abort}\) if no eligible proof has
been recorded, whether because of an internal algebraic failure or because
computation remains incomplete.  This deadline event is not part of the
probability in Lemma~\ref{lem:good23-product}.

The verifier parses one canonical prefix, reconstructs the statement digest
and the entire transcript, checks every work predicate, derives the selected
18 queries, authenticates the five opening sections, verifies the terminal
and relation equations, and checks every fold transition.  It need not run
\(\mathsf{Good}_{23}\): that predicate governs the honest prover's privacy
and availability policy, not acceptance of adversarial proofs.

The state-changing instruction consumes these bytes from a finalized,
pre-uploaded proof account.  Proof-account creation, chunk upload, and
finalization precede the accepting instruction.  ``One transaction'' in this
work refers to complete verification together with the nullifier and pool
mutation, not to that setup lifecycle.
